
\documentclass[12pt,a4paper]{article}

\usepackage[T1]{fontenc}
\usepackage[utf8]{inputenc}
\usepackage[spanish,es-noshorthands]{babel}
\usepackage{lmodern}
\usepackage{microtype}
\usepackage{geometry}

\geometry{
    a4paper,
    left=3cm,
    right=3cm,
    top=3cm,
    bottom=3cm
}

\title{De la virtud del liberalismo}
\author{Vicente Domínguez Arca}
\date{}

\begin{document}

\maketitle

Quizás no se necesiten muchos ratos de imaginación para llegar a querer que la \emph{areté} aristotélica vino a dar sentido humanista o moderno a la \emph{virtus} latina. Sin embargo, por lo que sea, la llamamos virtud hoy en día, que, aunque guarde como gran secreto la \emph{areté} y su anexa capacidad de justicia, sigue usando las letras que apelan a aquel coraje bélico. Y, a mí, que no me suelen convencer las casualidades, la virtud del liberalismo me parece más una \emph{virtus} que una \emph{areté}. No cerraré el texto sin presentar mi argumento, y pretendo que sea demasiado sólido.

En las sociedades occidentales de hoy en día da la impresión de que hay instituciones que no es que sean intocables, sino que semejan ser incontestables. Una de ellas es la propiedad privada. El adalid del liberalismo. El máximo exponente de lo que se considera la libertad de una persona. Y suena idílico, poético. El problema quizás no resida en enunciar el porqué de la propiedad privada, ni siquiera el qué; quizás podamos desplazar la discusión a cómo se mantiene la propiedad privada. Entiéndase: en una sociedad liberal, o neoliberal, un sujeto puede sentir su máxima expresión de libertad en la acción de poseer algo para su uso a discreción. Y realmente, si nos queremos poner poéticos, en su cara podemos ver ejemplos de aquello abstracto que llamamos felicidad o autorrealización. No preguntaré por más sentimientos. Pero sí me gustaría calmar los ánimos de ese sujeto anónimo preguntando cómo pretende continuar poseyendo a discreción la cosa en cuestión. Porque necesitará una estrategia.

Quizás pueda apelar directamente a que será la sociedad la que deba dar garantías de ello, y no deja de parecerme extraordinariamente curioso cómo un discurso liberal necesita tan pronto, en sus bases, apelar a algo social. Y, si levantamos suficientes ampollas, quizás obtengamos como respuesta que las garantías de protección las ejercerá el propio poseedor. Y ahí es donde aparece lo que considero la \emph{virtus}, o virtud, del liberalismo. Coraje bélico, valentía en el campo de batalla. Y, siendo sincero, diré que así, a mi parecer, todo encaja, y encaja bien. Pero a lo largo de la historia hemos visto demasiados ejemplos de valerosos caballeros con gran virtud que acabaron por necesitar vivir más de su pasado de lo que podían vivir con su presente. Porque sí, si algo le sucede a lo que apreciamos como cuerpo humano es que envejece y se debilita.

Sin embargo, a menudo, las propiedades para las que queremos garantías o que defendemos van mucho más allá en el tiempo que el periodo de virtud. Y así, inevitablemente, se vuelve imprescindible que sean otros, con otras virtudes, individuales o comunes, quienes defiendan la propiedad de uno, debilitado ya por el tiempo. Sin quererlo, se necesita echar mano de un entramado de derechos y deberes; acudir a la justicia acaba siendo lo más sensato para cualquier liberal que, sin poder en virtud, quiere defender su propiedad igualmente. Pero los derechos y los deberes son cuestiones sociales. Sí tiene sentido liberal referirse a la casa de una determinada persona, pero no tiene sentido decir que una persona tiene un determinado deber o un determinado derecho a su nombre, sino que esa persona puede encontrarse en disposición jurídica de ejercer un derecho o respetar un deber. Se trata de algo intrínsecamente social.

La dualidad que subyace en los argumentos liberales es algo paradójico, tanto que este autor se permite señalar un error de base argumental, una incoherencia incontestable, una acción meridianamente deliberada. Y eso es un problema, claro que lo es, porque al amparo de la virtud del liberalismo hoy en día se descarta \emph{de facto} cualquier otro tipo de propuesta de organización social; incluso se ridiculiza abiertamente o se transporta a escenarios históricos muy cuestionables. Desde aquí, desde mi singular posición libre, que no liberal, me gustaría decir que la virtud del liberalismo es un reflejo de su talón de Aquiles. Quizás sería interesante que nos planteáramos una organización social sin virtud o, más aún, por qué debemos vivir subyugados tácitamente a un sistema que crece al amparo de una virtud de cuestionable \emph{areté}.

\end{document}

